\documentclass[11pt]{article}
\usepackage[a4paper, top=18mm, bottom=18mm, left=20mm, right=20mm]{geometry}
\usepackage[T2A]{fontenc}
\usepackage[utf8]{inputenc}
\usepackage[ukrainian]{babel}
\usepackage{graphicx}
\setlength{\parindent}{0pt}
\setlength{\parskip}{4pt}
\pagestyle{empty}
\begin{document}
%begin
{\large\textbf{Контрольна робота}}

\textbf{Варіант \No\,1}\hfill \textit{ПІБ:}\,\underline{\hspace{60mm}}
\vspace{4mm}

%beginitem
1. Що буде виведено за виконання фрагменту коду?

%code
try:
    m = 4
    while m<10:
        m += 1
    print(m, end=';')
except:
    print('Z')

%code

%enditem
%beginitem
2. %goodpath:Scripts/2018/Mod2018_1/03-good.txt
Відмітити все, що є коректним оператором С++:
( 1 ) int 3a=2;

( 2 ) return 'a'<'c';

( 3 ) x=2'*'7;

( 4 ) j=i;

%enditem
%beginitem
3. Що буде виведено за виконання фрагменту коду?
Записати ВИРАЗ, значенням якого є множина, що складається з квадратівцілих чисел від 15 до 2005 (включно), що не діляться на 5.
%enditem
%beginitem
4. Що буде виведено за виконання фрагменту коду?

a,b,c=4,1,9
def f(a):
global c    a=4
    b=3
    c=2
    return a+b+c

a,b,c=0,8,7
print(f(b),a,b,c)

%enditem
%beginitem
5. Що буде виведено за виконання фрагменту коду?

\begin{center}
	\adjustimage{max size={0.8\linewidth}{0.8\paperheight}}
	{../BinTree/trees_exp/45.png}
\end{center}
Указати послідовність міток вершин дерева, зображеного на рис., що утвориться в результаті обходу дерева в прямому порядку. Мітки записувати через пробіл.\par Обхід дерева починається з кореня (на рис. це верхня вершина).
%answer:45 прямому

%enditem
%beginitem
6. Що буде виведено за виконання фрагменту коду?
Записати ВИРАЗ, значенням якого є список, що складається з квадратівцілих чисел від 15 до 2005 (включно), що не діляться на 5, записаних в обернено\-му порядку.

%enditem
%beginitem
7. Що буде виведено за виконання фрагменту коду?
def f(a,b):
      c=55
      if b==0: return 4
      if a<=-3: c=1
      else: c=9
      return c
   print(f(-4,-4))
%enditem
%beginitem
8. Що буде виведено за виконання фрагменту коду?

%code
e=[10,11,12,13]
e[-5:-5]='11'
print(e)
%code


%enditem
%beginitem
9. Що буде виведено за виконання фрагменту коду?
Задано тип вузла списку
struct Node \{int n; Node *prev, *next;\};
У списку зберігаються послідовні цілі числа від -25 до 49. Вказівник p встановлено на вузол, в якому зберігається число -4.
Чому дорівнює значення p->prev->prev->prev->prev->prev->n ?
%enditem
%beginitem
10. Що буде виведено за виконання фрагменту коду?

%code
print(6 != 6.0 > True <= 3)
%code

%enditem
%beginitem
11. Що буде виведено за виконання фрагменту коду?

print('R', end=' ')
m=10
while m>4:
    m+=-3
print(m)


%enditem
%beginitem
12. Що буде виведено за виконання фрагменту коду?

a,b,c=8,7,9
def f(a,b=6,c):
    print(a,b,c,end="")

f(a=4,1,a=5)
print(a,b,c)


%enditem
%beginitem
13. Що буде виведено за виконання фрагменту коду?

%code
a,b,c=4,1,9
def f(a):
global c    a=4
    b=3
    c=2
    return a+b+c

a,b,c=0,8,7
print(f(b),a,b,c)
%code

%enditem
%beginitem
14. Що буде виведено за виконання фрагменту коду?
Написати програму, що запитує у користувача значення дійсних параметрів $x$ та $y$, обчислює для них значення виразу 
$$\frac{\sqrt x + 27}{(\arccos x + 0.6) (y + 73)}$$ 
та повiдомляє введенi данi та результат обчислення.


%enditem
%beginitem
15. Що буде виведено за виконання фрагменту коду?
a = 4
b = 1
c = 9

def f(a):
global c    a = 4
    b += 3
    c = 2
    return a + b + c

a = 0
b = 8
c = 7

try:
    print(f(b), a, b, c, sep=';')
except:
    print('Z', a, b, c, sep=';')

%enditem
%beginitem
16. Що буде виведено за виконання фрагменту коду?
%code

def f(n):
    print("f", end="");
    return n

print(f(-4) and f(7))
%code

%enditem
%beginitem
17. Що буде виведено за виконання фрагменту коду?
try:
    m = 4
    while m<10:
        m += 1
    print(m, end=';')
except:
    print('Z')

%enditem
%beginitem
18. Що буде виведено за виконання фрагменту коду?
"6.0"==3.0
%enditem
%beginitem
19. Що буде виведено за виконання фрагменту коду?
not6==6.0 and notTrue<=3 and 7!=3.0
%enditem
%beginitem
20. Що буде виведено за виконання фрагменту коду?

%code
e=[10,11,12,13]
e.extend([1,2])
print(e)
%code


%enditem
%beginitem
21. Що буде виведено за виконання фрагменту коду?

print('R', end=' ')
e=('0','1','2','3','4','5','6','7','8','9')
print(e[-10:-10:2])


%enditem
%beginitem
22. Що буде виведено за виконання фрагменту коду?

%code
#include <iostream>

int f(int a, int b){
    int c = 55;
    if (b == 0)
        return 4;
    if (a <= -3)
         c = 1;
    else 
        c = 9;
    return c;
}

int main(){
    std::cout << f(-4, -4);
    return 0;
}
%code

%enditem
%beginitem
23. Що буде виведено за виконання фрагменту коду?
%code
for e in range(-10,-10,2):
    if e>7:
        continue
        print(e,end=' ')
    if e>7:
        break
else:
    print(e, end=' ')
print(e, end=' ')
%code

%enditem
%beginitem
24. Що буде виведено за виконання фрагменту коду?
Записати ВИРАЗ, значенням якого є список, що складається з цілих чисел від 15 до 2005 (включно), причому спочатку за зростанням записано всі, що не діляться на 5, а потім решту за зростанням.
%enditem
%beginitem
25. Що буде виведено за виконання фрагменту коду?

\begin{center}
	\adjustimage{max size={0.8\linewidth}{0.8\paperheight}}
	{../BinTree/trees_exp/45.png}
\end{center}
Указати послідовність міток вершин дерева, зображеного на рис., що утвориться в результаті обходу дерева в прямому порядку. Мітки записувати через пробіл.\par Алгоритм починає роботу з кореня (на рис. це верхня вершина).
%answer:45 прямому

%enditem
%beginitem
26. Що буде виведено за виконання фрагменту коду?
%code
#include <iostream>
using namespace std;

int f(int &x, int &y){
    x = 6;
    y+= 3;
    return x;
}

int main(){
    int a = 2, b = 9;
    a = f(a, a);
    cout << a << ":" << b <<':';
    {
        int a = 1, b = 4;
        cout << ((a>6) && ((b+=1) > 6)) << ':';
        cout << a << ":" << b << ':';
    }
    {
        inta = 8;
        cout << a << ':';
    }
    cout << a << ":" << b << endl;
    return 0;
}
%code


%enditem
%beginitem
27. Що буде виведено за виконання фрагменту коду?
try:
    for e in range(-10, -10, 2):
        if e>7:
            continue
        print(e, end=';')
        if e>7:
            break
    else:
        print(e,end=';')
    print(e)
except:
    print('ERR')

%enditem
%beginitem
28. Що буде виведено за виконання фрагменту коду?
def f(n):
    print('f', end='')
    return n

try:
    print(f(-4) and f(7))
except:
    print('ERR')

%enditem
%beginitem
29. Що буде виведено за виконання фрагменту коду?

%code
print(8/8//8*8)
%code

%enditem
%beginitem
30. Що буде виведено за виконання фрагменту коду?

%code
cout << (!6==6.0 && !true<= 3 && 7!=3.0);
%code

%enditem
%beginitem
31. Що буде виведено за виконання фрагменту коду?
e=('a','b','c','d','e',0,1,2,4); print(e[11])
%enditem
%beginitem
32. Що буде виведено за виконання фрагменту коду?
Задано тип вузла списку struct Node {int n; Node *prev, *next;};
У списку зберігаються послідовні цілі числа від -25 до 49. Вказівник p встановлено на вузол, в якому зберігається число -4.
Чому дорівнює значення p->prev->prev->prev->prev->prev->n ?
%enditem
%beginitem
33. Що буде виведено за виконання фрагменту коду?

%code
a = 4
b = 1
c = 9

def f(a):
global c    a=4
    b=3
    c=2
    return a+b+c

a, b, c = 0,8,7
try:
    print(f(b), a, b, c, sep=';')
except:
    print('Z', a, b, c, sep=';')
%code

%enditem
%beginitem
34. Що буде виведено за виконання фрагменту коду?

print('R', end=' ')
m=4
while m<10:
    m+=1
print(m)


%enditem
%beginitem
35. Що буде виведено за виконання фрагменту коду?

%code
try:
    m = 10
    while m>4:
        m += -3
    print(m, end=';')
except:
    print('Z')

%code

%enditem
%beginitem
36. Що буде виведено за виконання фрагменту коду?

%code
cout << (8 * 8 * 8 * 8);
%code

%enditem
%beginitem
37. Що буде виведено за виконання фрагменту коду?

print('R', end=' ')
m=10
while m>4:
    m+=-3
print(m)


%enditem
%beginitem
38. Що буде виведено за виконання фрагменту коду?

print('R', end=' ')
e=[10,11,12,13,14]
e[:-5]='13'
print(e)


%enditem
%beginitem
39. Що буде виведено за виконання фрагменту коду?
%code
if (4 == 4)
    cout << "x";
else
    cout << "b";
%code


%enditem
%beginitem
40. Що буде виведено за виконання фрагменту коду?
Рядки журналу через двокрапку містять ім'я користувача (ідентифікатор), час події,тип події, код події, наприклад
\begin{verbatim}
s="username : 2022-05-05 13:26 : ERROR : 404"
\end{verbatim}
у коді 
\begin{verbatim}
res = re.fullmatch('???', s) 
s = ''
code_str = ??? 
\end{verbatim}
чим треба замінити кожен з ???, щоб значенням code\_str був код події? 



%enditem
%beginitem
41. Що буде виведено за виконання фрагменту коду?

%code
a = 4
b = 1
c = 9

def f(a):
global c    a=4
    b+=3
    c=2
    return a+b+c

a, b, c = 0, 8, 7
try:
    print(f(b), a, b, c, sep=';')
except:
    print('Z', a, b, c, sep=';')
%code

%enditem
%beginitem
42. Що буде виведено за виконання фрагменту коду?
try:
    res = "6.0"==3.0j
    if isinstance(res, bool):
        print(f'bool;{res}')
    elif isinstance(res, int):
        print(f'int;{res}')
    elif isinstance(res, float):
        print(f'float;{res:.2f}')
    else:
        print('???')
except:
    print('ERR')

%enditem
%beginitem
43. %goodpath:Scripts/2019/Mod2019_1/Lex/03-good.txt
Відмітити все, що є коректним оператором С++:
( 1 ) if a<b: a=0 else: b=0

( 2 ) --n

( 3 ) if ('x'><'y'): a=6

( 4 ) j=i=13

%enditem
%beginitem
44. Що буде виведено за виконання фрагменту коду?
Знайти кількість відношень на n-елементній множині, які нерефлексивні та несиметричні
%enditem
%beginitem
45. Що буде виведено за виконання фрагменту коду?
Скільки 2017-елементних підмножин множини натуральних від 1 до 20172017 містять рівно одне число, що більше 172004
%enditem
%beginitem
46. Що буде виведено за виконання фрагменту коду?
try:
    print(4, end=';')
    print(int('c1'), end=';')
    print(9, end=';')
except ZeroDivisionError:
    print(0, end=';')
except TypeError:
    print(8, end=';')
except:
    print('ERR', end=';')
else:
    print(7, end=';')
finally:
    print(3)

%enditem
%beginitem
47. Що буде виведено за виконання фрагменту коду?

%code
print(1**1.0*True)
%code

%enditem
%beginitem
48. Що буде виведено за виконання фрагменту коду?
В умовах попередньої задачі було виконано p->prev=p->prev->prev;

Чому дорівнює значення p->prev->prev->prev->prev->n ?
%enditem
%beginitem
49. Що буде виведено за виконання фрагменту коду?

print('R', end=' ')
m=4
while m<10:
    m+=1
print(m)


%enditem
%beginitem
50. Що буде виведено за виконання фрагменту коду?

def f(n):
    print("f", end="");
    return n

print(f(-4) and f(7))

%enditem
%beginitem
51. Що буде виведено за виконання фрагменту коду?
Записати ВИРАЗ, значенням якого є множина, що складається з цілих чисел від 15 до 2005 (включно), що не діляться на 5.
%enditem
%beginitem
52. Що буде виведено за виконання фрагменту коду?

%code
print('R', end=' ')
e=[10,11,12,13,14]
e[:-5]='13'
print(e)
%code


%enditem
%beginitem
53. Що буде виведено за виконання фрагменту коду?
Рядки журналу через = містять ім'я користувача (ідентифікатор), час події,тип події, код події, наприклад
\begin{verbatim}
s="username = 2022-05-05 13:26 = ERROR = 404"
\end{verbatim}
у коді 
\begin{verbatim}
res = re.fullmatch('???', s) 
s = ''
code_str = ??? 
\end{verbatim}
чим треба замінити кожен з ???, щоб значенням code\_str був код події? 



%enditem
%beginitem
54. Що буде виведено за виконання фрагменту коду?

%code
print("6.0"==3.0)
%code

%enditem
%beginitem
55. Що буде виведено за виконання фрагменту коду?

%code
e = ('a','b','c','d','e',0,1,2,4)
print(e[11])
%code


%enditem
%beginitem
56. Що буде виведено за виконання фрагменту коду?

%code
#include <iostream>
using namespace std;

int a = 4, b = 1, c = 9;

int f(int &a){
int c;    a = 4;
    b -= 3;
    c = 2;
    return a + b + c;
}

int main(){
    inta = 0;
    intb = 8;
    intc = 7;
    cout << f(b) << ':';
    cout << a << ':' << b << ':' << c;
    return 0; 
}
%code

%enditem
%beginitem
57. Що буде виведено за виконання фрагменту коду?

print('R', end=' ')
e=[10,11,12,13,14]
e.extend([1,2])
print(e)


%enditem
%beginitem
58. %goodpath:Scripts/2019/Mod2019_1/Lex/01-good.txt
Відмітити все, що є однією правильною лексемою:
( 1 ) "abc'

( 2 ) 0x1a

( 3 ) x+23

( 4 ) "a\na"

%enditem
%beginitem
59. Що буде виведено за виконання фрагменту коду?
Змінні \kw{next} та \kw{prev} вузла списку позначають наступний та попередній за ним вузол відповідно. Починаючи з голови, у список записано послідовні цілі числа від~$-25$ до~$49$. Нехай змінна \kw{p} позначає вузол, що зберігає значення~$-4$.
Яке число зберігається у вузлі \kw{p.prev.prev.prev.prev.prev} ?

%enditem
%beginitem
60. Що буде виведено за виконання фрагменту коду?

%code
print(6!=6.0!=True!=3)
%code

%enditem
%beginitem
61. Що буде виведено за виконання фрагменту коду?

print('R', end=' ')
e=('0','1','2','3','4','5','6','7','8','9')
print(e[-10:-10:2])


%enditem
%beginitem
62. Що буде виведено за виконання фрагменту коду?

%code
print('R', end=' ')
for e in range(-5,-5,2):
    if e>7:
        continue
        print(e, end=' ')
    if e>7:
        break
    else:
        print(e, end=' ')
    print(e, end=' ')
%code

%enditem
%beginitem
63. %goodpath:Scripts/2018/Mod2018_1/01-good.txt
Відмітити все, що є однією правильною лексемою:
( 1 ) 0,5

( 2 ) false

( 3 ) x+23

( 4 ) x23

%enditem
%beginitem
64. Що буде виведено за виконання фрагменту коду?
Записати ВИРАЗ, значенням якого є множина, що складається з цілих чисел від 15 до 2005 (включно), що не діляться на 5.
%enditem
%beginitem
65. Що буде виведено за виконання фрагменту коду?
for e in range(-10,-10,2):
    if e>7:
        continue
        print(e,end=' ')
    if e>7:
        break
else:
    print(e, end=' ')
print(e, end=' ')

%enditem
%beginitem
66. Що буде виведено за виконання фрагменту коду?

a,b,c=4,1,9
def f(a):
global c    a=4
    b+=3
    c=2
    return a+b+c

a,b,c=0,8,7
print(f(b),a,b,c)

%enditem
%beginitem
67. Що буде виведено за виконання фрагменту коду?
try:
    res = 8/8//8*8
    if isinstance(res, bool):
        print(f'bool;{res}')
    elif isinstance(res, int):
        print(f'int;{res}')
    elif isinstance(res, float):
        print(f'float;{res:.2f}')
    else:
        print('???')
except:
    print('Z')

%enditem
%beginitem
68. Що буде виведено за виконання фрагменту коду?
Записати ВИРАЗ, значенням якого є словник, ключами якого є цілі числа від 15 до 2005 (включно), що не діляться на 5, а ключам відповідають їх квадрати 

%enditem
%beginitem
69. Що буде виведено за виконання фрагменту коду?
Записати ВИРАЗ, значенням якого є список, що складається з квадратівцілих чисел від 15 до 2005 (включно), що не діляться на 5, записаних в оберненому порядку.
%enditem
%beginitem
70. Що буде виведено за виконання фрагменту коду?

%code
for e in range(-5,-5+5,2):
    if e>7:
        break
    print(e, end=' ')
    e=-5
    if e>7:
        break
else:
    print(e, end=' ')
print(e, end=' ')
%code

%enditem
%beginitem
71. Що буде виведено за виконання фрагменту коду?
8/8//8*8
%enditem
%beginitem
72. Що буде виведено за виконання фрагменту коду?
try:
    res = 1**1.0*True
    if isinstance(res, bool):
        print(f'bool;{res}')
    elif isinstance(res, int):
        print(f'int;{res}')
    elif isinstance(res, float):
        print(f'float;{res:.2f}')
    else:
        print('???')
except:
    print('ERR')

%enditem
%beginitem
73. Що буде виведено за виконання фрагменту коду?

%code
print('R', end=' ')
e = ('a','b','c','d','e',0,1,2,4)
print(e[11])
%code


%enditem
%beginitem
74. Що буде виведено за виконання фрагменту коду?

%code
"6.0"==3.0
%code

%enditem
%beginitem
75. Що буде виведено за виконання фрагменту коду?

%code
#include <iostream>
using namespace std;

bool f(int n){
    cout<<"f";
    return n;
}

int main(){
    cout<<(f(-4) && f(7));
    return 0;
}
%code

%enditem
%beginitem
76. Що буде виведено за виконання фрагменту коду?

%code
#include <iostream>

int f(int c){
    int y = 55;
    if (c == 0) 
        return 4;
    if (c <= -3)
         y = 1;
    else
         y = 9;
    return y;
}

int main(){
    std::cout << f(-4);
    return 0;
}
%code


%enditem
%beginitem
77. Що буде виведено за виконання фрагменту коду?

\begin{center}
	\adjustimage{max size={0.8\linewidth}{0.8\paperheight}}
	{../BinTree/trees_exp/45.png}
\end{center}
Указати послідовність міток вершин дерева, зображеного на рис., що утвориться в результаті обходу дерева в прямому порядку. ̳тки записувати через пробіл.\par Обхід дерева починається з кореня (на рис. це верхня вершина).
%answer:45 прямому
%enditem
%beginitem
78. Що буде виведено за виконання фрагменту коду?

print('R', end=' ')
for e in range(-10,-10,2):
    if e>7:
        continue
        print(e, end=' ')
    if e>7:
        break
    else:
        print(e, end=' ')
    print(e)

%enditem
%beginitem
79. Що буде виведено за виконання фрагменту коду?
Записати ВИРАЗ, значенням якого є множина, що складається з квадратівцілих чисел від 15 до 2005 (включно), що не діляться на 5.

%enditem
%beginitem
80. Що буде виведено за виконання фрагменту коду?
a,b,c=4,1,9
   def f(a):
     global c     a=4
     b=3
     c=2
     return a+b+c
   a,b,c=0,8,7
   print(f(b),a,b,c)
%enditem
%beginitem
81. Що буде виведено за виконання фрагменту коду?
for e in range(-10,-10,2):
    if e>7:
        continue
        print(e,end=' ')
    if e>7:
        break
    else:
        print(e,end=' ')
    print(e)

%enditem
%beginitem
82. Що буде виведено за виконання фрагменту коду?
%Оригинал был утерян. Требует отладки!
%code
_c = 0

class C:
    _c = 1
    
    def __init__(self):
        self._c = 2
        _c = 3
        C._c = 4

obj = C()
obj._c = 5
try:
    print(_c, end=' ')
    print(obj._c, end=' ')
    print(C._c, end=' ')
    print(obj._c, end=' ')
    print(obj._C__c)
except:
    print('ERR')
%code

%enditem
%beginitem
83. Що буде виведено за виконання фрагменту коду?

\begin{center}
	\adjustimage{max size={0.8\linewidth}{0.8\paperheight}}
	{../15BinTree/trees_exp/45.png}
\end{center}
Указати послідовність міток вершин дерева, зображеного на рис., що утвориться в результаті обходу дерева в прямому порядку. Мітки записувати через пробіл.\par Алгоритм починає роботу з кореня (на рис. це верхня вершина).
%answer:45 прямому

%enditem
%beginitem
84. Що буде виведено за виконання фрагменту коду?
$a$ є цілочисловим масивом типу $numpy.ndarray$ розмірів $4{\times}4$. Сума елементів масиву $a$ дорівнює 25. Чому дорівнює сума елементів масиву $a$ після виконання 
\begin{verbatim}
a += 25
\end{verbatim}


Якщо код некоректний, то в якості відповіді зазначити ERROR.



%enditem
%beginitem
85. Що буде виведено за виконання фрагменту коду?
e=[10,11,12,13,14]; e[:-5]='13'; print(e)
%enditem
%beginitem
86. Що буде виведено за виконання фрагменту коду?
Рядки журналу через = містять ім'я користувача (ідентифікатор), час події,тип події, код події, наприклад
\begin{verbatim}
s="username = 2022-05-05 13:26 = ERROR = 404"
\end{verbatim}
у коді 
\begin{verbatim}
res = re.fullmatch('???', s) 
s = ''
code_str = ??? 
\end{verbatim}
чим треба замінити кожен з ???, щоб значенням code\_str був код події? 



%enditem
%beginitem
87. Що буде виведено за виконання фрагменту коду?
1**1.0*True
%enditem
%beginitem
88. Що буде виведено за виконання фрагменту коду?

%code
#include <iostream>
using namespace std;

int f(int a, int b){
    int c = 55;
    if (b == 0)
        return 4;
    if (a <= -3)
         c = 1;
    else 
        c = 9;
    return c;
}

int main(){
    cout << f(-4, -4);
    return 0;
}
%code

%enditem
%beginitem
89. Що буде виведено за виконання фрагменту коду?
Написати програму, що запитує у користувача значення дійсних параметрів $x$ та $y$, обчислює для них значення виразу та повiдомляє введенi данi та результат обчислення.
$$\frac{\sqrt x + 27}{(\arccos x + 0.6) (y + 73)}$$ 


%enditem
%beginitem
90. Що буде виведено за виконання фрагменту коду?
Записати ВИРАЗ, значенням якого є список, що складається з цілих чисел від 15 до 2005 (включно), причому спочатку за зростанням записано всі, що не діляться на 5, а потім решту за зростанням.

%enditem
%beginitem
91. Що буде виведено за виконання фрагменту коду?

%code
a,b,c=4,1,9
def f(a):
    a=4
    b+=3
    c=2
    return a+b+c

a,b,c=0,8,7
print(f(b),a,b,c)
%code

%enditem
%beginitem
92. Що буде виведено за виконання фрагменту коду?
"6.0"==3.0j
%enditem
%beginitem
93. Що буде виведено за виконання фрагменту коду?
e=[10,11,12,13,14]; e[:-5]='13'; print(e)
%enditem
%beginitem
94. Що буде виведено за виконання фрагменту коду?
def f(c):
      y=55
      if c==0: return 4
      if c<=-3: y=1
      else: y=9
      return y
   print(f(-4))
%enditem
%beginitem
95. Що буде виведено за виконання фрагменту коду?
try:
    m = 10
    while m>4:
        m += -3
    print(m, end=';')
except:
    print('Z')

%enditem
%beginitem
96. Що буде виведено за виконання фрагменту коду?

%code
cout << (6 == 6.0 <= true != 3);
%code

%enditem
%beginitem
97. Що буде виведено за виконання фрагменту коду?
m=10
   while m>4: m+=-3
   print(m)
%enditem
%beginitem
98. Що буде виведено за виконання фрагменту коду?

%code
cout << 9 * 9 * 9 * 9;
%code

%enditem
%beginitem
99. Що буде виведено за виконання фрагменту коду?
e=[10,11,12,13,14]; e.extend([1,2]); print(e)
%enditem
%beginitem
100. Що буде виведено за виконання фрагменту коду?
Записати ВИРАЗ, значенням якого є список, що складається з цілих чисел від 15 до 2005 (включно), причому спочатку за зростанням записано всі, що не діляться на 5, а потім решту за зростанням.
%enditem




































































































%end
\newpage
\end{document}


